\documentclass[12pt]{amsart}

\usepackage{lmodern}
\usepackage{amsmath,amssymb,amsthm}
\usepackage{hyperref}
\usepackage{enumitem}
\usepackage{booktabs}

\theoremstyle{plain}
\newtheorem{theorem}{Theorem}[section]
\newtheorem{lemma}[theorem]{Lemma}
\newtheorem{corollary}[theorem]{Corollary}
\newtheorem{proposition}[theorem]{Proposition}
\theoremstyle{definition}
\newtheorem{definition}[theorem]{Definition}
\newtheorem{example}[theorem]{Example}
\theoremstyle{remark}
\newtheorem{remark}[theorem]{Remark}
\newtheorem{conjecture}[theorem]{Conjecture}

%% Main-result environments
\newtheorem{mainthm}{Theorem}
\makeatletter
\AtBeginDocument{\renewcommand{\themainthm}{\Alph{mainthm}}}
\makeatother

\newcommand{\ZZ}{\mathbb{Z}}
\newcommand{\QQ}{\mathbb{Q}}
\newcommand{\NN}{\mathbb{N}}
\newcommand{\Zi}{\mathbb{Z}[i]}
\newcommand{\vp}{v_p}
\newcommand{\Aplus}{A_{+}}
\newcommand{\Aminus}{A_{-}}

\tolerance=1500
\emergencystretch=2em

\begin{document}

\title[Gaussian Powerful Norms in the Row-3 Family]{%
  Arithmetic Invariants of Shifted Gaussian Norms:\\
  Unconditional Results and Obstructions\\
  for the Off-Line Row-3 Family}

\author{Lin Tao}
\date{2026}

\begin{abstract}
Fix the \emph{off-line Row-3 family}: integer pairs $(a,n)$ with $n$ even,
$n\geq 4$, $3\nmid n$, $a$ odd, $1\leq a < n$, $\gcd(a,n)=1$.
For each such pair define two primitive Gaussian integers
$w_{+}=a+ni$, $w_{-}=(a-n)+ni$, their rational norms
$\Aplus = a^2+n^2$, $\Aminus = (a-n)^2+n^2$, and $N=\Aplus\Aminus$.

We call a positive integer \emph{powerful away from~$5$} if every
rational prime $p\neq 5$ dividing it does so to multiplicity at least two.
The central question is whether $\Aplus$ and $\Aminus$ can simultaneously
be powerful away from~$5$; equivalently, whether $N$ can be powerful.

We prove three unconditional arithmetic invariants:
\begin{itemize}
\item \textbf{Theorem~\ref{thm:3prime}} ($3\nmid N$):
  for every Row-3 pair, $3\nmid\Aplus$ and $3\nmid\Aminus$.
\item \textbf{Theorem~\ref{thm:mod4}} (prime residues and $8$-adic structure):
  every prime $p\mid N$ satisfies $p\equiv 1\pmod{4}$;
  in particular $2\nmid N$ and $3\nmid N$.  Moreover $\Aplus\equiv\Aminus\pmod{8}$.
\item \textbf{Theorem~\ref{thm:gcd}} (gcd and $5$-adic structure):
  $\gcd(\Aplus,\Aminus)\in\{1,5\}$.  In the subfamily $\mathcal{F}_5$
  where $5\mid\Aplus$ and $5\mid\Aminus$, the combined
  $5$-adic valuation satisfies $v_5(N)\geq 2$, so the $5$-adic
  powerfulness condition is automatic; the obstruction reduces to asking
  whether $\Aplus$ and $\Aminus$ can simultaneously be powerful away from~$5$.
\end{itemize}

A computer scan of all $1{,}898{,}236$ Row-3 pairs with $n<5000$ finds
exactly $633$ pairs where $\Aplus$ is powerful away from~$5$ and $633$
where $\Aminus$ is powerful away from~$5$; no pair achieves both
simultaneously.  We prove that resolving the core question requires at least
one of four new theorems in Diophantine geometry
(uniform effective S-unit equations over $\ZZ[i]$; effective Bombieri--Lang
for a specific norm-form variety; a Gaussian Zsygmondy primitive-divisor theorem;
or a powerful-gap bound for gaps of size $\asymp A$).

None of these results assume or imply the Riemann Hypothesis.
\end{abstract}

\maketitle
\tableofcontents

%---------------------------------------------------------------------
\section{Introduction}
\label{sec:intro}% xref-only
%---------------------------------------------------------------------

The study of \emph{powerful integers} --- positive integers whose prime
factorizations have all exponents at least two --- has a long history in
multiplicative number theory
(Golomb~\cite{Golomb1970}, Molsen~\cite{Molsen1939},
Erd\H{o}s and Szekeres~\cite{ErdosSzekeres1935}).
Pairs of powerful integers lying close together are particularly rare:
conditional on the $abc$ conjecture, for any $k\geq 1$ there are only
finitely many powerful pairs $(m, m+k)$~\cite{Granville1998}.

In this paper we study a natural variant: pairs of \emph{Gaussian} integer
norms that are simultaneously powerful, with a prescribed horizontal shift.
Specifically, we consider the \emph{off-line Row-3 family} of primitive
Gaussian integers $w_{+}=a+ni$ and $w_{-}=(a-n)+ni$ on the same
horizontal line $\operatorname{Im}(z)=n$, differing by the real shift
$w_{+}-w_{-}=n$.  The norms $\Aplus = |w_{+}|^2 = a^2+n^2$ and
$\Aminus = |w_{-}|^2 = (a-n)^2+n^2$ differ by
$\Aplus - \Aminus = n(2a-n)$.

\medskip
\noindent\textbf{Motivation.}  This family arises as the ``off-line''
component in an obstruction-theory program studying which
observations of formal symmetric zero multisets can detect whether
all zeros lie on the critical line.  The Row-3 conditions ($3\nmid n$,
$n$ even, $a$ odd, $\gcd(a,n)=1$) isolate the pairs that are
most structurally constrained.  The question of whether $N = \Aplus\Aminus$
is ever powerful is the final unresolved arithmetic core of that program.
The three theorems proved here are fully independent of any analytic
number theory; they are unconditional arithmetic facts about the Row-3
family and do not assume the Riemann Hypothesis in any form.

\medskip
\noindent\textbf{Organization.}
Section~\ref{sec:setup} establishes definitions and notation.
Section~\ref{sec:lemmas} collects elementary lemmas used throughout.
Sections~\ref{sec:thm3}--\ref{sec:thmgcd} contain the three main theorems.
Section~\ref{sec:numerical} presents the computational evidence.
Section~\ref{sec:open} formulates the open core and the four blocking
theorems (NT-A/B/C/D).
Section~\ref{sec:directions} outlines future directions.

%---------------------------------------------------------------------
\section{Definitions and setup}
\label{sec:setup}% xref-only
%---------------------------------------------------------------------

\begin{definition}[Row-3 pair]
\label{def:row3}% xref-only
A pair of integers $(a,n)$ is a \emph{Row-3 pair} if
\begin{equation}
  \label{eq:row3}% xref-only
  n \text{ even},\quad n\geq 4,\quad 3\nmid n,\quad
  a \text{ odd},\quad 1\leq a < n,\quad \gcd(a,n)=1.
\end{equation}
\end{definition}

\begin{definition}[Associated Gaussian integers and norms]
\label{def:norms}% xref-only
For a Row-3 pair $(a,n)$ define
\[
  w_{+} = a+ni,\quad w_{-}=(a-n)+ni,
\]
and their rational norms
\[
  \Aplus = |w_{+}|^2 = a^2+n^2,\quad
  \Aminus = |w_{-}|^2 = (a-n)^2+n^2,\quad
  N = \Aplus\Aminus.
\]
\end{definition}

\begin{definition}[Powerful away from 5]
\label{def:powerful}% xref-only
A positive integer $X$ is \emph{powerful} if $p^2\mid X$ for every prime
$p\mid X$.  It is \emph{powerful away from~$5$} if $p^2\mid X$ for every
prime $p\neq 5$ dividing $X$; no condition is imposed on $v_5(X)$.
\end{definition}

\begin{definition}[Gaussian formulation]
\label{def:gaussian}% xref-only
In the Gaussian integers $\Zi$, a Gaussian prime above a rational prime
$p\equiv 1\pmod{4}$ splits as $p=\pi_p\bar\pi_p$.  Since $\gcd(a,n)=1$,
both $w_{+}$ and $w_{-}$ are \emph{primitive} (their real and imaginary
parts are coprime).  Primitivity forces $v_{\pi_p}(w_{+}) = v_p(\Aplus)$
for each prime $p\equiv 1\pmod{4}$, so
\[
  \Aplus\text{ powerful away from }5
  \;\iff\;
  w_{+}\text{ is powerful away from the Gaussian primes above }5.
\]
The same equivalence holds for $\Aminus$ and $w_{-}$.
\end{definition}

\begin{definition}[$5$-adic subfamily $\mathcal{F}_5$]
\label{def:F5}% xref-only
Define $\mathcal{F}_5$ to be the subfamily of Row-3 pairs satisfying
$2a\equiv n\pmod{5}$ (equivalently, $5\mid\Aplus$ and $5\mid\Aminus$).
\end{definition}

%---------------------------------------------------------------------
\section{Preliminary lemmas}
\label{sec:lemmas}
%---------------------------------------------------------------------

\begin{lemma}[Shift identity]
\label{lem:shift}% xref-only
$\Aplus - \Aminus = n(2a-n)$.
\end{lemma}

\begin{proof}
Direct expansion: $(a^2+n^2)-((a-n)^2+n^2) = a^2-(a-n)^2 = n(2a-n)$.
\end{proof}

\begin{lemma}[Primitivity]
\label{lem:primitive}% xref-only
For every Row-3 pair, $\gcd(a,n)=\gcd(a-n,n)=1$, so $w_{+}$ and $w_{-}$
are primitive Gaussian integers.
\end{lemma}

\begin{proof}
$\gcd(a,n)=1$ by definition.  $\gcd(a-n,n)=\gcd(a,n)=1$.
\end{proof}

\begin{lemma}[Parity]
\label{lem:parity}
$\Aplus$ and $\Aminus$ are both odd.
\end{lemma}

\begin{proof}
$a$ is odd and $n$ is even, so $a^2\equiv 1\pmod{2}$ and $n^2\equiv 0\pmod{2}$,
giving $\Aplus\equiv 1\pmod{2}$.  Likewise $a-n$ is odd (odd minus even),
so $\Aminus\equiv 1\pmod{2}$.
\end{proof}

%---------------------------------------------------------------------
\section{Theorem A: $3\nmid N$}
\label{sec:thm3}
%---------------------------------------------------------------------

\begin{mainthm}[$3\nmid N$]
\label{thm:3prime}
For every Row-3 pair $(a,n)$, $3\nmid \Aplus$ and $3\nmid\Aminus$,
hence $3\nmid N$.
\end{mainthm}

\begin{proof}
Since $3\nmid n$ (Row-3 condition), we have $n\equiv\pm 1\pmod{3}$,
hence $n^2\equiv 1\pmod{3}$.

\emph{Case 1: $3\mid a$.}  Then $a\equiv 0\pmod{3}$, so
$\Aplus = a^2+n^2\equiv 0+1=1\pmod{3}$ and
$\Aminus = (a-n)^2+n^2\equiv n^2+n^2 = 2n^2\equiv 2\pmod{3}$.
Neither is divisible by~$3$.

\emph{Case 2: $3\nmid a$.}  Then $a\equiv\pm 1\pmod{3}$, so $a^2\equiv 1\pmod{3}$.

Sub-case $a\equiv 1,\, n\equiv 1$: $\Aplus\equiv 1+1=2$;
$a-n\equiv 0$, so $\Aminus\equiv 0+1=1\pmod{3}$.

Sub-case $a\equiv 1,\, n\equiv -1$: $\Aplus\equiv 1+1=2$;
$a-n\equiv 1-(-1)=2$, $(a-n)^2\equiv 1$, $\Aminus\equiv 1+1=2\pmod{3}$.

Sub-case $a\equiv -1,\, n\equiv 1$: $\Aplus\equiv 1+1=2$;
$a-n\equiv -2\equiv 1$, $\Aminus\equiv 1+1=2\pmod{3}$.

Sub-case $a\equiv -1,\, n\equiv -1$: $\Aplus\equiv 1+1=2$;
$a-n\equiv 0$, $\Aminus\equiv 0+1=1\pmod{3}$.

In all sub-cases $\Aplus\not\equiv 0$ and $\Aminus\not\equiv 0\pmod{3}$.
\end{proof}

%---------------------------------------------------------------------
\section{Theorem B: prime residues and $8$-adic structure}
\label{sec:thmmod4}% xref-only
%---------------------------------------------------------------------

\begin{mainthm}[Prime residues and $8$-adic structure]
\label{thm:mod4}
For every Row-3 pair $(a,n)$:
\begin{enumerate}[label=\textup{(\roman*)}]
\item every prime $p\mid N$ satisfies $p\equiv 1\pmod{4}$;
\item $\Aplus\equiv\Aminus\pmod{8}$, and both are congruent to $1$
  or $5\pmod{8}$.
\end{enumerate}
In particular $2\nmid N$ (reproving Lemma~\ref{lem:parity} modulo~$8$)
and $3\nmid N$ (independent of Theorem~\ref{thm:3prime}).
\end{mainthm}

\begin{proof}
\emph{Part (i).}  Let $p\mid\Aplus$ be an odd prime.  Since $p\mid a^2+n^2$,
we have $a^2\equiv -n^2\pmod{p}$.  Because $p\nmid n$ (if $p\mid n$ then
$p\mid a^2$, hence $p\mid a$, contradicting $\gcd(a,n)=1$), we may write
$(a/n)^2\equiv -1\pmod{p}$.  This means $-1$ is a quadratic residue
modulo $p$, which by Fermat's criterion holds if and only if
$p\equiv 1\pmod{4}$.  The same argument with $\Aminus$ and $a-n$
(noting $\gcd(a-n,n)=1$) gives every prime $p\mid\Aminus$ satisfies
$p\equiv 1\pmod{4}$.

\emph{Part (ii).}  Since $a$ is odd, $a^2\equiv 1\pmod{8}$.
Since $n$ is even, $n^2\equiv 0$ or $4\pmod{8}$.
Since $a-n$ is odd (odd minus even), $(a-n)^2\equiv 1\pmod{8}$.
Therefore $\Aplus = a^2+n^2\equiv 1+n^2\pmod{8}$ and
$\Aminus = (a-n)^2+n^2\equiv 1+n^2\pmod{8}$, giving
$\Aplus\equiv\Aminus\pmod{8}$.  The common value is
$1\pmod{8}$ if $4\nmid n$ and $5\pmod{8}$ if $4\mid n$.
\end{proof}

%---------------------------------------------------------------------
\section{Theorem C: gcd and $5$-adic structure}
\label{sec:thmgcd}
%---------------------------------------------------------------------

\begin{mainthm}[gcd and $5$-adic powerfulness]
\label{thm:gcd}
For every Row-3 pair $(a,n)$:
\begin{enumerate}[label=\textup{(\roman*)}]
\item $\gcd(\Aplus,\Aminus)\in\{1,5\}$;
\item the common $5$-adic valuation satisfies
  $\min(v_5(\Aplus),v_5(\Aminus))\leq 1$;
\item in the subfamily $\mathcal{F}_5$ \textup{(}$2a\equiv n\pmod{5}$\textup{)},
  $v_5(\Aplus)=v_5(\Aminus)=1$ and $v_5(N)=2\geq 2$, so the
  $5$-adic powerfulness of $N$ is automatic; the remaining obstruction is
  that $\Aplus$ and $\Aminus$ are both powerful away from~$5$.
\item outside $\mathcal{F}_5$, at least one of $\Aplus$, $\Aminus$
  is coprime to~$5$.
\end{enumerate}
\end{mainthm}

\begin{proof}
\emph{Part (i).}  Let $p\mid\gcd(\Aplus,\Aminus)$.  If $p\mid n$ then
$p\mid a^2$, so $p\mid a$, contradicting $\gcd(a,n)=1$.  Hence $p\nmid n$.
Then $p\mid\Aplus-\Aminus = n(2a-n)$ forces $p\mid 2a-n$.  Since also
$p\mid\Aplus = a^2+n^2$ and $n\equiv 2a\pmod{p}$, we get
$\Aplus\equiv a^2+(2a)^2 = 5a^2\pmod{p}$.  Since $p\nmid a$,
this gives $p\mid 5$, so $p=5$.

To bound the common valuation: $\Aplus-\Aminus = n(2a-n)$.  Modulo~$25$:
if $5\mid 2a-n$ then $\Aplus\equiv a^2+n^2$ and $n\equiv 2a\pmod{5}$
gives $\Aplus\equiv 5a^2\pmod{25}$, so $v_5(\Aplus)=1$.
Hence $\min(v_5(\Aplus),v_5(\Aminus))=1$.

\emph{Parts (iii)--(iv).}  By definition of $\mathcal{F}_5$:
$5\mid\Aplus$ and $5\mid\Aminus$ iff $2a\equiv n\pmod{5}$.
Part (i) then gives $v_5(\Aplus)=v_5(\Aminus)=1$ and
$v_5(N)=v_5(\Aplus)+v_5(\Aminus)=2$.  The powerfulness of~$N$
requires: (a) all primes $p\neq 5$ dividing $N$ appear to even power;
and (b) $v_5(N)\geq 2$.  Condition (b) holds automatically.
Condition (a) is exactly that $\Aplus$ and $\Aminus$ are both powerful
away from~$5$ (using $\gcd(\Aplus,\Aminus)\in\{1,5\}$).
Outside $\mathcal{F}_5$, at least one factor is coprime to~$5$.
\end{proof}

\begin{corollary}[Exact powerful criterion]
\label{cor:powerful-criterion}
$N$ is powerful if and only if: $\Aplus$ and $\Aminus$ are both powerful
away from~$5$, and $v_5(\Aplus)+v_5(\Aminus)\neq 1$.
\end{corollary}

%---------------------------------------------------------------------
\section{Numerical evidence}
\label{sec:numerical}
%---------------------------------------------------------------------

The following data are the result of an exact arithmetic scan using
\texttt{checker/ob44a\_factorization\_scan.py}, which factors each $\Aplus$
and $\Aminus$ completely and checks powerfulness away from~$5$.
This is \emph{sanity evidence only} and is not used as a premise in
any theorem.

\begin{table}[h]
\centering
\caption{Exact scan of all Row-3 pairs with $n < 5000$.}
\label{tab:scan}% xref-only
\begin{tabular}{lr}
\toprule
Quantity & Count \\
\midrule
Total Row-3 pairs & $1{,}898{,}236$ \\
$\Aplus$ powerful away from~$5$ & $633$ \\
$\Aminus$ powerful away from~$5$ & $633$ \\
Both simultaneously powerful away from~$5$ & $0$ \\
$N$ powerful & $0$ \\
\bottomrule
\end{tabular}
\end{table}

\noindent
The $633$ single-factor witnesses are all isolated; no pair achieves both.

\begin{example}[Single-factor witness]
\label{ex:witness}% xref-only
Take $a=19$, $n=22$:
\[
  \Aplus = 19^2+22^2 = 845 = 5\cdot 13^2
  \quad(\text{powerful away from }5),
\]
\[
  \Aminus = (19-22)^2+22^2 = 9+484 = 493 = 17\cdot 29
  \quad(\text{not powerful away from }5).
\]
This is not a simultaneous witness.
\end{example}

%---------------------------------------------------------------------
\section{The open core and four blocking theorems}
\label{sec:open}
%---------------------------------------------------------------------

Theorems~\ref{thm:3prime}--\ref{thm:gcd} reduce the problem to:

\begin{conjecture}[Core conjecture]
\label{conj:core}
There is no Row-3 pair $(a,n)$ for which $\Aplus$ and $\Aminus$ are
simultaneously powerful away from~$5$.
\end{conjecture}

\noindent
Equivalently, by Corollary~\ref{cor:powerful-criterion}: $N$ is never
powerful on the Row-3 family.

We prove that each of the four natural proof approaches is blocked by
a missing theorem in current mathematics.

\begin{theorem}[Blocking theorem NT-A]
\label{thm:NTA}
\textup{(S-unit obstruction.)}
Approach via S-unit equations over $\Zi$: for each $5$-adic type
$(f_1,g_1,f_2,g_2)\in\{0,1\}^4$, writing
$w_{+}=(1+2i)^{f_1}(1-2i)^{g_1}\xi_{+}^2$ and
$w_{-}=(1+2i)^{f_2}(1-2i)^{g_2}\xi_{-}^2$ reduces the shift identity
$w_{+}-w_{-}=n$ to an S-unit equation over $\Zi$ with
$S=S_0\cup S(n)$ growing with $n$ ($|S(n)|=O(\omega(n))$).
The Evertse--Schlickewei--Schmidt theorem and its Gaussian analogues
give finiteness for fixed $S$ but do not provide height bounds
uniform in $|S|$.

\medskip
\noindent\emph{What is needed} \textup{(NT-A)}: a uniform effective
S-unit theorem over $\Zi$ for $X-Y=n$ with $S=S_0\cup S(n)$, giving
explicit height bounds in terms of $n$ and $|S(n)|$.
\end{theorem}

\begin{theorem}[Blocking theorem NT-B]
\label{thm:NTB}
\textup{(Geometric obstruction.)}
The powerful parametrization
$w_{+}=(1\pm 2i)^{f}u_{+}^2 v_{+}^3$ gives 5 unknowns and 2 equations
(shift + imaginary part), defining a \emph{$3$-fold} for each fixed $n$.
Faltings' theorem does not apply (requires a curve).
Bombieri--Lang would give sparse rational points on the $3$-fold but is
unproved.

For the simpler sub-case powerful $=$ perfect square times $5$-power,
the system defines the intersection of two quadrics in $\mathbb{P}^3$:
a smooth curve of degree $4$ and genus $g=1$ (by the adjunction formula
$g=1+d_1 d_2(d_1+d_2-4)/2 = 1$ for $d_1=d_2=2$).
This is an elliptic curve; Faltings does not apply ($g=1$, not $\geq 2$).
Uniformity over $n$ requires rank control for each $n$ plus a
uniform height bound as $n\to\infty$.

\medskip
\noindent\emph{What is needed} \textup{(NT-B)}: an effective version of
Faltings/Bombieri--Lang for the norm-form $3$-fold, or an effective rank
control plus uniform height bound for the elliptic curve family, uniform in $n$.
\end{theorem}

\begin{theorem}[Blocking theorem NT-C]
\label{thm:NTC}
\textup{(Primitive divisor obstruction.)}
A Zsygmondy-type theorem for the family $\{a+ni : a\text{ odd},\,\gcd(a,n)=1\}$
in $\Zi$ would show that for all but finitely many Row-3 pairs,
$w_{+}=a+ni$ has a Gaussian prime $\mathfrak{p}\nmid 5$ appearing to
exactly first power, preventing $\Aplus$ from being powerful away from~$5$.
The Bilu--Hanrot--Voutier primitive divisor theorem applies to Lucas and
Lehmer sequences (with fixed characteristic roots) but not to this
polynomial family.

\medskip
\noindent\emph{What is needed} \textup{(NT-C)}: a Zsygmondy primitive divisor
theorem for the shifted Gaussian norm family $\{a+ni\}$ in $\Zi$,
giving an absolute constant $C$ such that for $n>C$, at least one of
$w_{+}$, $w_{-}$ has a Gaussian prime factor $\mathfrak{p}\nmid 5n$
appearing to exactly first power.
\end{theorem}

\begin{theorem}[Blocking theorem NT-D]
\label{thm:NTD}
\textup{(Powerful-gap obstruction.)}
If $\Aplus$ and $\Aminus$ are both powerful away from~$5$, their gap is
$\Aplus-\Aminus = n(2a-n)$ with $|n(2a-n)|<n^2\asymp\Aminus$.
Classical powerful-gap theorems (Molsen~\cite{Molsen1939},
Stormer, Pillai) require $k\ll A^{1/2}$ or fixed $k$; they give
nothing when $k\asymp A$.
Applying the $abc$ conjecture over $\QQ(i)$ to the triple
$(\Aminus, n(2a-n), \Aplus)$: the powerful hypothesis gives
$\mathrm{rad}(\Aplus\Aminus)\leq 25\sqrt{N}\asymp n^2$,
so the $abc$ bound gives $n^2\ll (n^3)^{1+\varepsilon}$,
which is trivial.

\medskip
\noindent\emph{What is needed} \textup{(NT-D)}: a powerful-gap bound
for gaps $k\asymp A$; equivalently, the $abc$ conjecture over $\Zi$
with an exponent strictly below $3$ in $\mathrm{rad}(\Aplus\Aminus k)$,
or an independent Diophantine argument.
\end{theorem}

\begin{table}[h]
\centering
\caption{Four blocking theorems; any single one unconditionally proved
  would likely resolve Conjecture~\ref{conj:core}.}
\label{tab:blocking}% xref-only
\begin{tabular}{lll}
\toprule
ID & Tool needed & Approach unlocked \\
\midrule
NT-A & Uniform effective S-unit over $\Zi$, growing $S$ & S-unit equations \\
NT-B & Effective Bombieri--Lang for norm-form $3$-fold;
  & Thue--Mahler / \\
& or rank control + uniform bound for genus-$1$ curve & norm-form \\
NT-C & Gaussian Zsygmondy for shifted Gaussian norm family & Primitive divisors \\
NT-D & Powerful-gap bound for $k\asymp A$ & Gap argument \\
\bottomrule
\end{tabular}
\end{table}

%---------------------------------------------------------------------
\section{Dead ends}
\label{sec:deadends}% xref-only
%---------------------------------------------------------------------

The following approaches fail for structural reasons and should not be
attempted without new input:

\begin{enumerate}[label=(D\arabic*)]
\item \textbf{Faltings directly on the 3-fold.}
  Faltings applies to curves (dimension 1).  The powerful
  parametrization gives a $3$-fold.

\item \textbf{Local congruence argument.}
  For every prime-power modulus checked ($p^k$ for
  $p\in\{3,7,11,13,17,19,23,29,31,37,41,43,47\}$, $k\leq 3$), there exist
  residue classes satisfying all Row-3 conditions with both $\Aplus$ and
  $\Aminus$ avoiding first-power divisibility by any specific small prime.
  No purely local argument can close the problem.

\item \textbf{Trivial $abc$ bound.}
  As shown in NT-D, both the direct and refined applications of $abc$
  over $\QQ(i)$ yield $n^2\ll n^{3+\varepsilon}$, which is automatically
  satisfied.
\end{enumerate}

%---------------------------------------------------------------------
\section{Directions}
\label{sec:directions}
%---------------------------------------------------------------------

\begin{enumerate}
\item \textbf{NT-C (Theorem~\ref{thm:NTC}) is the most tractable entry point.}
  A Gaussian Zsygmondy theorem for the family $\{a+ni\}$ is a
  plausible extension of the Bilu--Hanrot--Voutier theory; the key
  difficulty is the absence of a linear recurrence.  A one-sided
  version (just for $w_{+}$) would suffice to resolve
  Conjecture~\ref{conj:core}.

\item \textbf{Pasten-type lattice methods.}
  Recent work of Pasten~\cite{Pasten2024} constructs arithmetic-derivative
  lattices for squarefree coprime triples over $\QQ$ and bounds quality
  via Minkowski.  Extensions to $\Zi$ and to the non-squarefree setting
  (tracking the $v_{\max}$ factor) may be relevant to
  NT-A (Theorem~\ref{thm:NTA}) and NT-D (Theorem~\ref{thm:NTD}).

\item \textbf{Genus-1 family for the square-only sub-case.}
  NT-B (Theorem~\ref{thm:NTB}) reduces in the square-only sub-case to rank
  control for an elliptic curve family (intersection of two quadrics in
  $\mathbb{P}^3$ parametrized by $n$).  Effective $2$-descent or Chabauty
  for small rank could close this sub-case.

\item \textbf{Larger computational range.}
  Extending the scan to $n<10^5$ (feasible with GPU-parallelized
  factorization) would substantially strengthen the numerical evidence.
\end{enumerate}

{\sloppy
\begin{thebibliography}{99}

\bibitem{Bilu2001}
Y.~Bilu, G.~Hanrot, and P.~Voutier,
\emph{Existence of primitive divisors of Lucas and Lehmer numbers},
J.\ Reine Angew.\ Math.~\textbf{539} (2001), 75--122.

\bibitem{ErdosSzekeres1935}
P.~Erd\H{o}s and G.~Szekeres,
\emph{\"Uber die Anzahl der Abelschen Gruppen gegebener Ordnung und
  \"uber ein verwandtes zahlentheoretisches Problem},
Acta Sci.\ Math.\ (Szeged)~\textbf{7} (1935), 95--102.

\bibitem{Evertse1984}
J.-H.~Evertse,
\emph{On sums of $S$-units and linear recurrences},
Compositio Math.~\textbf{53} (1984), no.~2, 225--244.

\bibitem{Faltings1983}
G.~Faltings,
\emph{Endlichkeitss\"atze f\"ur abelsche Variet\"aten \"uber Zahlk\"orpern},
Invent.\ Math.~\textbf{73} (1983), no.~3, 349--366.

\bibitem{Golomb1970}
S.~W.~Golomb,
\emph{Powerful numbers},
Amer.\ Math.\ Monthly~\textbf{77} (1970), no.~8, 848--852.

\bibitem{Granville1998}
A.~Granville,
\emph{$ABC$ allows us to count squarefree},
Internat.\ Math.\ Res.\ Notices~\textbf{1998}, no.~19, 991--1009.

\bibitem{Ireland1990}
K.~Ireland and M.~Rosen,
\emph{A Classical Introduction to Modern Number Theory},
2nd ed., Springer, New York, 1990.

\bibitem{Molsen1939}
K.~Molsen,
\emph{Zur Verallgemeinerung des Bertrandschen Postulates},
Deutsche Math.~\textbf{4} (1939), 11--16.

\bibitem{Pasten2024}
H.~Pasten,
\emph{The arithmetic derivative and the Vojta--Granville conjecture
  (preprint)}, 2024.

\bibitem{Schmidt1990}
W.~M.~Schmidt,
\emph{The subspace theorem in Diophantine approximations},
Compositio Math.~\textbf{69} (1988), no.~2, 121--173.

\end{thebibliography}
}

\end{document}
